\section{Problem Formulation and Analysis}
\label{sec:formulation}

We formulate multi-turn on-policy distillation in an interactive environment.
At each OPD update, let \(\pi_{\theta_{\mathrm{old}}}\) denote the current
student policy used for data collection, and let \(\pi_\theta\) denote the
student policy being optimized. For each input \(x\sim\mathcal D\), the agent
interacts with an environment over \(T\) decision steps.

At interaction step \(t\), the agent observes an interaction history
\(H_t=(O_1,A_1,O_2,A_2,\ldots,A_{t-1},O_t)\), where \(O_t\) is the current
observation and \(A_s\) is the action taken at step \(s\). Given \(H_t\), a policy
\(\pi\) chooses an action \(A_t\sim\pi(\cdot\mid x,H_t)\), and the environment returns
the next observation \(O_{t+1}\sim\mathcal E(\cdot\mid x,H_t,A_t)\). Thus a step is an
interaction step consisting of a policy action followed by an environment response,
rather than merely a token-generation step.

For any policy \(\pi\), the policy-environment system induces a step-\(t\) history
occupancy distribution
\(d_{\pi,\mathcal E}^{t}(h_t\mid x)=\Pr_{\pi,\mathcal E}(H_t=h_t\mid x)\), which
factorizes recursively: if \(h_{t+1}=(h_t,a_t,o_{t+1})\), then
\[
d_{\pi,\mathcal E}^{t+1}(h_{t+1}\mid x)
=
d_{\pi,\mathcal E}^{t}(h_t\mid x)\,
\pi(a_t\mid x,h_t)\,
\mathcal E(o_{t+1}\mid x,h_t,a_t).
\]
In particular, the current student \(\pi_{\theta_{\mathrm{old}}}\) induces the
student interaction occupancy
\(d_{\theta_{\mathrm{old}}}^{t}\equiv d_{\pi_{\theta_{\mathrm{old}}},\mathcal E}^{t}\),
while the teacher \(\pi_T\) induces the teacher interaction occupancy
\(d_T^{t}\equiv d_{\pi_T,\mathcal E}^{t}\). Once a history \(h_t\) is collected, the
teacher is queried for a target action distribution \(\pi_T(\cdot\mid x,h_t)\).
Therefore, in interactive OPD the design choice is not only which target
distribution to distill from, but also which interaction histories to query the
teacher on.

\paragraph{Ideal interactive improvement objective.}

Let \(q_t^\star(\cdot\mid x,h_t)\) denote the ideal improvement target at interaction
step \(t\) and history \((x,h_t)\). The ideal objective evaluates the updated student
\(\pi_\theta\) on histories that the current student \(\pi_{\theta_{\mathrm{old}}}\)
will actually encounter under the environment interaction:
\begin{equation}
\label{eq:ideal}
\mathcal R^\star(\theta;\theta_{\mathrm{old}})
=
\mathbb E_{x\sim\mathcal D}
\left[
\sum_{t=1}^{T}
\alpha_t\,
\mathbb E_{H_t\sim d_{\theta_{\mathrm{old}}}^{t}(\cdot\mid x)}
\left[
\ell\left(
\pi_\theta(\cdot\mid x,H_t),
q_t^\star(\cdot\mid x,H_t)
\right)
\right]
\right],
\end{equation}
where \(\alpha_t\ge 0\) controls the relative importance of different
interaction steps; unless a schedule is specified we take \(\alpha_t\equiv 1\)
(uniform across steps), so that ReOPD's reliability signal -- including its step-decay --
is carried entirely by the weight \(w_t\) (Section~\ref{sec:aopd}), not by
\(\alpha_t\). This objective captures student relevance: the student should
improve on histories generated by its own interaction with the environment.

In practice, \(q_t^\star\) is unavailable. OPD therefore replaces it with the
teacher distribution \(\pi_T(\cdot\mid x,h_t)\). However, the teacher is now
queried on histories generated by a particular roll-in process. If those
histories are far from the teacher's reliable interaction support, the teacher's
action distribution may not be a trustworthy improvement signal.

\paragraph{From online interaction to offline reuse.}

The occupancies \(d_{\pi,\mathcal E}^{t}\) and the ideal objective
\(\mathcal R^\star\) are defined through environment interaction, but they serve
only as \emph{conceptual targets}. Realizing them online would require, at every
update, rolling the current student through the environment and querying the
teacher afresh at each visited history -- precisely the per-step environment and
teacher cost we wish to avoid. We instead train off-environment from a fixed pool
of pre-collected teacher trajectories: at each recorded history \(h_t\) the
student proposes an action and the teacher conditional
\(\pi_T(\cdot\mid x,h_t)\) supplies the target, while no environment query is
needed because the prefix and observations are replayed from the trace.
Consequently the roll-in is the teacher pool, so the collected histories
follow \(\mathcal P_t\approx d_T^{t}\) rather than the student occupancy
\(d_{\theta_{\mathrm{old}}}^{t}\). The student occupancy that the ideal objective
targets is therefore \emph{never sampled}; obtaining it is exactly what
online interaction would buy. This is the crux of the offline setting: we must
approximate an objective defined over \(d_{\theta_{\mathrm{old}}}^{t}\) using
only samples from \(d_T^{t}\), which is what the reweighting below accomplishes.

\paragraph{Collected and effective history distributions.}

Write \(\mathcal P_t(\cdot\mid x)\) for the roll-in distribution histories are collected
from; in the offline setting it is the fixed teacher pool, so
\(\mathcal P_t\approx d_T^{t}\). A general weighted OPD objective is
\begin{equation}
\label{eq:weighted}
\mathcal L_{\mathcal P,w}(\theta;\theta_{\mathrm{old}})
=
\mathbb E_{x\sim\mathcal D}
\left[
\sum_{t=1}^{T}
\alpha_t\,
\mathbb E_{H_t\sim \mathcal P_t(\cdot\mid x)}
\left[
w_t(x,H_t)\,
\ell\left(
\pi_\theta(\cdot\mid x,H_t),
\pi_T(\cdot\mid x,H_t)
\right)
\right]
\right],
\end{equation}
where the weights are normalized per \((x,t)\), i.e.\
\(\mathbb E_{H_t\sim \mathcal P_t(\cdot\mid x)}[w_t(x,H_t)]=1\), so that \(w_t\) reshapes
\(\mathcal P_t\) into the per-step distribution \(\rho_t\); the position-only schedule of
Section~\ref{sec:aopd}, constant across histories at each \(t\), is instead normalized
across positions, as detailed there. Then \(w_t\) induces
an effective interaction-prefix distribution
\(\rho_t(h_t\mid x)=w_t(x,h_t)\,\mathcal P_t(h_t\mid x)\), and the objective becomes
\begin{equation}
\label{eq:rho-obj}
\mathcal L_{\rho}(\theta;\theta_{\mathrm{old}})
=
\mathbb E_{x\sim\mathcal D}
\left[
\sum_{t=1}^{T}
\alpha_t\,
\mathbb E_{H_t\sim \rho_t(\cdot\mid x)}
\left[
\ell\left(
\pi_\theta(\cdot\mid x,H_t),
\pi_T(\cdot\mid x,H_t)
\right)
\right]
\right].
\end{equation}
Thus the central design problem is not the raw weight \(w_t\), but the effective
history distribution \(\rho_t\).

\paragraph{Two sources of mismatch.}

The practical objective differs from the ideal interactive improvement objective
for two reasons. First, the effective history distribution \(\rho_t\) may differ
from the student-environment occupancy \(d_{\theta_{\mathrm{old}}}^{t}\).
Second, even when histories are relevant to the student, the teacher may be
unreliable on histories far from its own interaction support.

Define the step-wise teacher reliability error
\[
\epsilon_{T,t}^{\theta}(x,h_t)
=
\left|
\ell\left(
\pi_\theta(\cdot\mid x,h_t),
q_t^\star(\cdot\mid x,h_t)
\right)
-
\ell\left(
\pi_\theta(\cdot\mid x,h_t),
\pi_T(\cdot\mid x,h_t)
\right)
\right|.
\]
\begin{assumption}[Bounded improvement loss]
\label{asm:bounded}
The loss against the ideal target is uniformly bounded: there exists a constant
\(B<\infty\) such that
\(\bigl|\ell\left(\pi_\theta(\cdot\mid x,h_t),q_t^\star(\cdot\mid x,h_t)\right)\bigr|\le B\)
for all \((x,h_t,t)\).
\end{assumption}

Assumption~\ref{asm:bounded} is mild and covers the losses we use. It holds
automatically for any bounded divergence -- e.g.\ total variation (\(B=1\)) or
Jensen--Shannon (\(B=\log 2\)). For the per-token KL used in
Section~\ref{sec:aopd}, it holds whenever the target conditional is bounded below on
its support, \(q_t^\star(a\mid x,h_t)\ge p_{\min}>0\), since then
\(D_{\mathrm{KL}}(\pi_\theta\,\|\,q_t^\star)\le\log(1/p_{\min})=:B\); a softmax with
bounded logits (equivalently a temperature or label-smoothing floor) guarantees such
a \(p_{\min}\). Only the ideal-target loss needs to be bounded: the teacher-target
loss enters the bound solely through the difference \(\epsilon_{T,t}^{\theta}\).

\begin{proposition}[Two-sided decomposition]
\label{prop:bound}
Recall the ideal objective
\(\mathcal R^\star(\theta;\theta_{\mathrm{old}})\)~(\plaineqref{eq:ideal}), which scores
\(\pi_\theta\) against the ideal target \(q_t^\star\) on the student occupancy
\(d_{\theta_{\mathrm{old}}}^{t}\), and the replayed objective
\(\mathcal L_\rho(\theta;\theta_{\mathrm{old}})\)~(\plaineqref{eq:rho-obj}), which scores
\(\pi_\theta\) against the teacher \(\pi_T\) on the effective prefix distribution
\(\rho_t=w_t\,\mathcal P_t\). Under Assumption~\ref{asm:bounded} (with bound \(B\)),
\begin{equation}
\label{eq:bound}
\left|
\mathcal R^\star(\theta;\theta_{\mathrm{old}})
-
\mathcal L_\rho(\theta;\theta_{\mathrm{old}})
\right|
\le
\mathbb E_{x\sim\mathcal D}
\left[
\sum_{t=1}^{T}
\alpha_t
\left\{
2B\,
\mathrm{TV}
\left(
d_{\theta_{\mathrm{old}}}^{t}(\cdot\mid x),
\rho_t(\cdot\mid x)
\right)
+
\mathbb E_{H_t\sim \rho_t(\cdot\mid x)}
\left[
\epsilon_{T,t}^{\theta}(x,H_t)
\right]
\right\}
\right].
\end{equation}
Here \(\mathrm{TV}\) is the total-variation distance, \(\alpha_t\ge0\) the step
coefficients of~(\plaineqref{eq:ideal}), and
\(\epsilon_{T,t}^{\theta}(x,h_t)=\bigl|\ell(\pi_\theta,q_t^\star)-\ell(\pi_\theta,\pi_T)\bigr|\)
the step-wise teacher-reliability error, the gap between the ideal- and teacher-target
losses at \((x,h_t)\).
\end{proposition}

\begin{proof}
Fix \(x\) and \(t\), and abbreviate the ideal-target and teacher-target losses by
\(f_t(h)=\ell(\pi_\theta(\cdot\mid x,h),q_t^\star(\cdot\mid x,h))\) and
\(g_t(h)=\ell(\pi_\theta(\cdot\mid x,h),\pi_T(\cdot\mid x,h))\), so that
\(\mathcal R^\star=\mathbb E_{x}\bigl[\sum_t\alpha_t\,\mathbb E_{d_{\theta_{\mathrm{old}}}^{t}}[f_t]\bigr]\)
and \(\mathcal L_\rho=\mathbb E_{x}\bigl[\sum_t\alpha_t\,\mathbb E_{\rho_t}[g_t]\bigr]\).
Adding and subtracting \(\mathbb E_{\rho_t}[f_t]\),
\[
\mathbb E_{d_{\theta_{\mathrm{old}}}^{t}}[f_t]-\mathbb E_{\rho_t}[g_t]
=
\underbrace{\bigl(\mathbb E_{d_{\theta_{\mathrm{old}}}^{t}}[f_t]-\mathbb E_{\rho_t}[f_t]\bigr)}_{\text{occupancy shift}}
+
\underbrace{\bigl(\mathbb E_{\rho_t}[f_t]-\mathbb E_{\rho_t}[g_t]\bigr)}_{\text{teacher reliability}}.
\]
For the occupancy-shift term, Assumption~\ref{asm:bounded} gives
\(\|f_t\|_\infty\le B\), and since
\(\sum_h\bigl|d_{\theta_{\mathrm{old}}}^{t}(h\mid x)-\rho_t(h\mid x)\bigr|=2\,\mathrm{TV}(d_{\theta_{\mathrm{old}}}^{t},\rho_t)\),
\[
\bigl|\mathbb E_{d_{\theta_{\mathrm{old}}}^{t}}[f_t]-\mathbb E_{\rho_t}[f_t]\bigr|
\le
\|f_t\|_\infty\sum_h\bigl|d_{\theta_{\mathrm{old}}}^{t}(h\mid x)-\rho_t(h\mid x)\bigr|
\le
2B\,\mathrm{TV}\bigl(d_{\theta_{\mathrm{old}}}^{t}(\cdot\mid x),\rho_t(\cdot\mid x)\bigr).
\]
For the teacher-reliability term, Jensen's inequality and the definition of
\(\epsilon_{T,t}^{\theta}\) give
\(\bigl|\mathbb E_{\rho_t}[f_t-g_t]\bigr|\le\mathbb E_{\rho_t}[|f_t-g_t|]=\mathbb E_{\rho_t}[\epsilon_{T,t}^{\theta}]\).
Applying the triangle inequality, multiplying by \(\alpha_t\ge0\), summing over \(t\),
and taking \(\mathbb E_{x\sim\mathcal D}\) establishes~(\plaineqref{eq:bound}).
\end{proof}
The first term is the student interaction-occupancy mismatch: it is small when the
effective training histories match the histories that the current student will
encounter through environment interaction. The second term is the teacher
reliability mismatch: it is small when the teacher is queried on histories where its
action distribution is a reliable proxy for the ideal improvement target.

Fully student-on-policy OPD (\(\rho_t=d_{\theta_{\mathrm{old}}}^{t}\) for all \(t\))
zeroes the occupancy term but not the reliability one: student-generated actions can
steer the environment into histories unlikely under the teacher's own interaction
process, especially at later steps where its conditional is a poorly calibrated
target. Teacher-forced histories (\(\rho_t=d_T^{t}\)) do the reverse -- reliable near
the teacher's own occupancy, but not where the student actually goes at test time.
Neither extreme is uniformly optimal, so the effective distribution should balance
student relevance against teacher reliability.

\paragraph{Reliability-aware interaction-prefix distribution.}

The preceding decomposition motivates choosing, for each interaction step, an
effective history distribution
\[
\rho_t^\star
\in
\arg\min_{\rho_t}
\left\{
\lambda_{\mathrm{stu},t}\,
D\left(
\rho_t(\cdot\mid x),
d_{\theta_{\mathrm{old}}}^{t}(\cdot\mid x)
\right)
+
\lambda_{\mathrm{tea},t}\,
\mathcal E_T(\rho_t;x)
\right\},
\]
where \(D\) measures mismatch to the current student's occupancy and
\(\mathcal E_T(\rho_t;x)\) measures expected teacher unreliability under \(\rho_t\).
Since the true reliability error is generally unobserved, we use teacher support as a
practical surrogate, giving the bridge objective
\begin{equation}
\label{eq:bridge-obj}
\rho_t^\star
\in
\arg\min_{\rho_t}
\left\{
\lambda_{\mathrm{stu},t}\,
D_{\mathrm{KL}}
\left(
\rho_t(\cdot\mid x)
\,\Vert\,
d_{\theta_{\mathrm{old}}}^{t}(\cdot\mid x)
\right)
+
\lambda_{\mathrm{tea},t}\,
D_{\mathrm{KL}}
\left(
\rho_t(\cdot\mid x)
\,\Vert\,
d_T^{t}(\cdot\mid x)
\right)
\right\}.
\end{equation}
When the supports overlap, the solution is the geometric bridge
\begin{equation}
\label{eq:bridge}
\rho_t^\star(h_t\mid x)
\propto
\left[
d_{\theta_{\mathrm{old}}}^{t}(h_t\mid x)
\right]^{\gamma_t}
\left[
d_T^{t}(h_t\mid x)
\right]^{1-\gamma_t},
\qquad
\gamma_t
=
\frac{\lambda_{\mathrm{stu},t}}{\lambda_{\mathrm{stu},t}+\lambda_{\mathrm{tea},t}}.
\end{equation}
When \(\gamma_t=1\) the bridge becomes fully student-on-policy; when \(\gamma_t=0\) it
becomes teacher-supported roll-in. Intermediate values select histories that are both
likely under the current student's environment interaction and not too far from the
teacher's reliable interaction support.

\paragraph{Two regimes.}

The bound~(\plaineqref{eq:bound}) sums two competing terms: the occupancy term is
minimized by \(\gamma_t\!\to\!1\) (fully student-on-policy) and the reliability term by
\(\gamma_t\!\to\!0\) (teacher-supported roll-in), so which one dominates sets the right
\(\gamma_t\). Crucially the balance is \emph{per step}: the teacher-reliability term
grows with depth, as roll-ins drift into histories the teacher rarely visits, so the
ideal \(\gamma_t\) \emph{decreases along the trajectory} -- student-relevant
(\(\gamma_t\!\to\!1\)) at early steps, teacher-supported (\(\gamma_t\!\to\!0\)) only at
the deep steps where the teacher signal has degraded. This yields two regimes. When the
teacher stays reliable on student-induced histories -- a capable teacher close to the
student, short horizons, or early steps -- the occupancy term dominates, \(\gamma_t\!\to\!1\)
is best, and ordinary student-on-policy OPD is already strong. When the teacher's signal
degrades off its own support -- a large teacher--student gap, long horizons, or later
steps -- lowering \(\gamma_t\) \emph{at those steps} reduces the gap even though the
resulting histories are \emph{less} student-on-policy, which is counterintuitive if one
reads OPD as merely ``make it on-policy''. Section~\ref{sec:aopd} approximates this
depth-dependent retreat without tuning \(\gamma_t\) directly -- shifting supervision
mass toward earlier steps rather than reshaping the roll-in within each step --
applying on the fixed teacher-pool roll-in a step-decay whose steepness grows with the
teacher--student gap.

\paragraph{From distribution to weight.}

When histories come from \(\mathcal P_t\), any effective distribution
\(\rho_t^\star\) is realized by the normalized importance weight
\(w_t^\star=\rho_t^\star/\mathcal P_t\) (with
\(\mathbb E_{H_t\sim\mathcal P_t}[w_t^\star]=1\)), which recovers the weighted
objective~(\plaineqref{eq:weighted}) and can be applied equivalently by reweighting or
by sampling prefixes in proportion to it. The weight is thus an implementation device,
not the conceptual core: the principle is to choose, at each step, an effective history
distribution that balances student relevance and teacher reliability. Because exact
density ratios are high-variance, the concrete step-decay schedule we use in
practice is developed in Section~\ref{sec:aopd}. After optimizing the weighted
objective the student is refreshed,
\(\theta_{\mathrm{new}}\leftarrow\arg\min_{\theta}\mathcal L_{\rho}(\theta;\theta_{\mathrm{old}})\),
and the next OPD iteration sets
\(\theta_{\mathrm{old}}\leftarrow\theta_{\mathrm{new}}\).

\paragraph{Interpretation.}
On-policy distillation is not automatically optimal merely because histories are
sampled from the student: the history distribution sets both where the student is
trained and where the teacher is asked to supervise, so when the teacher is unreliable
on student-induced histories a reliability-aware distribution close to both occupancies
is the safer target.
